\labsection{7}{NLP Evaluation: Airline Tweet Sentiment Classification}{sec:lab07-nlp-b}

\begin{labproject}{Build and evaluate a three-class sentiment classifier}
\textbf{Scenario.} An airline support team wants to classify customer tweets as negative, neutral, or positive.\\
\textbf{Goal.} Build a leakage-safe baseline and Multinomial Naive Bayes benchmark from the single supplied \texttt{Tweets.csv} file.

\begin{labmode}
Use \texttt{all\_datasets/lab07-Tweets.csv} and run \texttt{all\_experiments/lab07\_airline\_sentiment\_worked.ipynb}. Save evidence under \texttt{results/lab07/}.
\end{labmode}

\begin{mlsteps}
  \item \textbf{Inspect and clean.} Keep tweet ID, text, sentiment, and airline. Report class balance, missing text, and duplicate text; remove exact duplicate text before splitting.
  \item \textbf{Create a fixed split.} Use seed 42 for a stratified 70/15/15 train/validation/test split and save tweet IDs with split labels.
  \item \textbf{Fit text features on training data only.} Normalize URLs/mentions and build unigram or unigram+bigrams with \texttt{CountVectorizer}.
  \item \textbf{Establish a baseline.} Predict the most frequent training class for all validation/test tweets.
  \item \textbf{Tune Naive Bayes on validation data.} Select n-gram range and smoothing parameter without inspecting test performance.
  \item \textbf{Evaluate once on test data.} Report accuracy, macro precision, macro recall, macro F1, and a 3-by-3 confusion matrix.
  \item \textbf{Inspect subgroup performance.} Report macro F1 per airline and review 12 deterministic test errors with structured error categories.
\end{mlsteps}

\begin{mltable}
\begin{tabularx}{\linewidth}{@{}>{\bfseries\leavevmode\color{MLNavy}}p{0.25\linewidth}XXXX@{}}
\toprule
\rowcolor{MLPaleBlue!92}
Method & Accuracy & Macro Precision & Macro Recall & Macro F1 \\
\midrule
Majority baseline & & & & \\
Multinomial NB & & & & \\
\bottomrule
\end{tabularx}
\end{mltable}

\requiredoutput{Executed notebook, split CSV, metrics CSV, confusion matrix, per-airline F1 CSV, and error-analysis CSV.}
\end{labproject}

\begin{verification}
A clean Python run with seed 42 must reproduce the split and test metrics. Test text must not influence vocabulary fitting, n-gram selection, smoothing selection, or any model choice.
\end{verification}

\begin{labdeliverable}
Push the notebook and results; create tag \texttt{lab07-final}. The README must compare baseline and Naive Bayes macro F1, describe the strongest/weakest airline subsets without treating them as airline-quality scores, and summarize the most common error types.
\end{labdeliverable}
